\section{Introduction}
\label{sec:intro}

Multi-angle total internal reflection fluorescence (MA-TIRF) microscopy images a sample at
several illumination incidence angles. Each angle probes a different, exponentially decaying
depth range near the coverslip, so the set of angle images encodes the axial ($z$) distribution
of the fluorophores. Recovering that distribution --- a 3D volume from a handful of 2D angle
images --- is a linear inverse problem, and a severely ill-posed one: the forward operator
determines only two to three depth directions out of ${\sim}50$ (Sec.~\ref{sec:problem}).

This report compares six reconstruction algorithms on that problem, and on 2D deconvolution as
a well-posed companion where the same algorithms are easy to sanity-check. It is written for
whoever maintains \texttt{matirf\_python\_plus} next: the goal is to know, for each algorithm,
\emph{how to use it}, \emph{when it works}, and \emph{what it cannot do}.

\paragraph{How to read it.}
Sec.~\ref{sec:problem} states the problem and the one fact that governs everything (the null
space). Sec.~\ref{sec:algorithms} recalls the six algorithms and their parameters.
Sec.~\ref{sec:setup} is the protocol (datasets, grids, metrics). Sec.~\ref{sec:results} is the
core --- the atlas of what wins where and with which parameters. Sec.~\ref{sec:limits} states
the intrinsic limits, and Sec.~\ref{sec:decision} is a one-page decision guide. In a hurry, read
only Sec.~\ref{sec:decision}.
